\subsection{Intrinsic fibration and even shifts}\label{app:known-periods}
The results cited as K in Sections 1 and 2 are known with the credit given there. The claim made for the text below concerns its re-derivation from the retained period charts.
\subsubsection{Re-derivation of FAM A.1.1}
\begin{proof}
\par\smallskip\noindent\textbf{Step 1.} We compute the line and its evaluation map, including the multiple fibres. A holomorphic function restricts to a constant on every connected compact reduced fibre. To descend functions on a neighbourhood we first use regular fibres; equality then extends from their dense union by the identity theorem on the smooth total space. Away from the multiple fibres, evaluation on local sections descends it to a holomorphic base function. On a finite root cover this follows by evaluation on the zero section; invariance under $s\mapsto\zeta s$ leaves precisely the convergent series in $s^m$. At the cusp, a smooth point of the reduced fibre admits a local section, and constancy on each nearby compact fibre gives the same descent there by the identity theorem. Consequently $f_*\mathcal O_X=\mathcal O_B$.

\par\smallskip\noindent\textbf{Step 2.} Put $e_\wedge=d\zeta_1\wedge d\zeta_2$ in relative fibre coordinates. At the finite centre $j$ let $a_1=2,a_2=1$ and $\zeta_j=e^{-2\pi i/m_j}$. The canonical frame $ds_j\wedge e_\wedge$ has character $\chi_j(s)=\zeta_j\det R_j(s)$, with $\chi_j(0)=\zeta_j^{1+a_j}$. Take its holomorphic logarithm $\chi_j=\zeta_j^{1+a_j}\exp\psi_j$, where $\psi_j(0)=0$. The finite group identity says $\sum_{r=0}^{m_j-1}\psi_j(\zeta_j^r s)=0$. Writing $\psi_j=\sum c_ns^n$, we obtain $c_n=0$ for $m_j\mid n$. Therefore
\[
 \varphi_j(s)=\sum_{m_j\nmid n}\frac{c_n}{1-\zeta_j^n}s^n
\]
is convergent and satisfies $\varphi_j(s)-\varphi_j(\zeta_js)=\psi_j(s)$. The invariant coefficient module of the relative canonical frame
$e^{\varphi_j}(ds_j\wedge e_\wedge)\otimes(dt_j)^{-1}$ is exactly
$s_j^{m_j-1-a_j}\mathbb C\{s_j^{m_j}\}$. Notice that $(dt_j)^{-1}$ denotes a pulled-back line-bundle frame; no division by a vanishing differential is made. Since the total-space cover is etale, the evaluation divisor downstairs is $(m_j-1-a_j)S_j$, namely $0$ and $2S_2$.

\par\smallskip\noindent\textbf{Step 3.} At the cusp the ordinary form
\[
 \Omega=\frac{dx_1}{x_1}\wedge\frac{dx_2}{x_2}\wedge dt_c
\]
is nowhere zero. Indeed in a maximal unimodular cone chart, with coordinates $z_0,z_1,z_2$, the height-one condition gives $t_c=z_0z_1z_2$ and unimodularity gives $\Omega=\pm dz_0\wedge dz_1\wedge dz_2$. The same equality after inverting coordinates treats the other charts. Integral shears and holomorphic torus multipliers depending on $t_c$ preserve this form: all their extra differential terms contain $dt_c$ and vanish in the wedge. Thus the form descends to the cusp quotient and evaluation has no zero there. This calculation also shows that changing any collar by a fibre translation does not alter the canonical transition factors.

\par\smallskip\noindent\textbf{Step 4.} We next determine $\mathfrak K=f_*\omega_{X/B}$. On the simply connected period cover let $h_6^2=E_6(\tau)$, and put $\Theta=E_4(\tau)^2h_6/\Delta_{\rm mod}(\tau)$. The relation $j(\tau)=1728t$ gives local degrees $\tau-\rho\sim s_1$, $\tau-i\sim s_2^2$ and $q_\tau=e^{2\pi i\tau}\sim t_c/1728$. The identities $E_4^3=1728t\Delta_{\rm mod}$ and $E_6^2=1728(t-1)\Delta_{\rm mod}$, together with the nonvanishing product $\Delta_{\rm mod}=q_\tau\prod_{n\ge1}(1-q_\tau^n)^{24}$ on the upper half-plane, give orders $2,1,-1$ for $\Theta$. The same identities exclude every other zero or pole. The square root exists because the zeros of $E_6(\tau)$ have even order and the cover is simply connected. Its signs are determined locally: at $\rho$, $\rho^3=-1$ and $h_6\ne0$, whereas at $i$, $h_6$ has order one in $s_2$. The weight transformation laws consequently give
\[
 \Theta(g_1z)=-\Theta(z)/\tau(z),\qquad
 \Theta(g_2z)=\Theta(z)/\tau(z),\qquad \Theta(g_0z)=\Theta(z).
\]
The invariant relative form $\Theta^{-1}e_\wedge$ is a unit in the finite local generator just calculated and has a simple zero in the cusp generator $e_\wedge=(2\pi i)^{-2}\Omega\otimes(dt_c)^{-1}$. It defines an isomorphism $\mathfrak K\simeq\mathcal O_B(p_0)$. The evaluation divisor now yields
\[
 \omega_X\simeq f^*\mathcal O_B(p_0)\otimes f^*\omega_B\otimes\mathcal O_X(2S_2).
\]
In the fixed coordinate $t$, the function $(t-1)^{-1}$ has divisor $p_0-p_2$, while $dt/(t-1)^2$ has divisor $-2p_2$. These explicit maps identify the last display with $\mathcal O_X(-4S_2+2S_2)$.

\par\smallskip\noindent\textbf{Step 5.} A meromorphic representative of a section of $\mathcal O_X(kS_2)$ is holomorphic off $S_2$, so the constancy argument makes it a pullback there. On the root cover its value along the zero section is meromorphic, with order at least $-k$, and is invariant under $s\mapsto-is$. Its Laurent series has only powers $s^{4r}$. The necessary and sufficient inequality is $4\operatorname{ord}_{p_2}b+k\ge0$, which proves the direct-image formula, including negative $k$.

\par\smallskip\noindent\textbf{Step 6.} Choose the canonical divisor section $s_{\rm can}$ of $\mathcal O_X(2S_2)$ and identify this line with $\omega_X^{-1}$ using the preceding maps. Then $s_{\rm can}^2=f^*\sigma_{p_2}$. Choose $\ell\in H^0(B,\mathcal O_B(p_2))$ with $\ell(p_2)\ne0$, and let $t_{\rm can}=f^*\ell$ in degree two. In degree $m=2r$ the monomials $s_{\rm can}^{2r-2j}t_{\rm can}^j$, $0\le j\le r$, are the pullbacks of the usual basis $\sigma_{p_2}^{r-j}\ell^j$. In degree $2r+1$, multiplication by $s_{\rm can}$ gives the same basis because the allowed base pole order remains $r$. These are all sections by the direct-image formula. Thus $U\mapsto s_{\rm can},V\mapsto t_{\rm can}$ is a graded algebra isomorphism, not only a Hilbert-series equality. The two degree-two sections have no common zero and give $f$ in the projective coordinate $[\sigma_{p_2}:\ell]$.

\par\smallskip\noindent\textbf{Step 7.} A biholomorphism pulls back the canonical line, its sections and their evaluations; it therefore induces $\alpha\in\operatorname{PGL}_2(\mathbb C)$ with $f'F=\alpha f$. The cusp fibre is reduced with nonnormal points. The two finite fibres have smooth reductions and nilpotence orders exactly three and four in their local equations $s^3=0,s^4=0$. These are intrinsic local-ring properties of the full fibres and distinguish all three. Every other fibre is a smooth reduced torus. Consequently $\alpha$ fixes $\infty,0,1$ individually, and the equation of a fractional linear transformation gives $\alpha=1$.
\end{proof}
\begin{proposition}[FAM A.1.2; T]\label{fam:A.1.2}
The normalized base change $N_j\times_{D_j}\Delta_j\to\Delta_j$, $t_j=s_j^{m_j}$, is the original smooth torus family $J'_j$. The maps $E_1\times F_1\to S_1$ and $E_2\times F_2\to S_2$ described below are finite etale covers of degrees $9$ and $8$.
\end{proposition}
\uses{Definition~\ref{def:conventions}; FAM A.1.1 and preceding appendix calculations.}
\begin{proof}
\par\smallskip\noindent\textbf{Step 1.} On $J'_j$, the deck generator is
$(s,\zeta)\mapsto(\zeta_js,R_j(s)\zeta+\Pi_c(g_js)v_j/m_j)$. It has exact order $m_j$. Its $r$th power changes the invariant circle coordinate $\gamma$ by $r\gamma(v_j)/m_j$; this is nonintegral for $0<r<m_j$, so the action is free. The quotient map is therefore a local biholomorphism. Its local base function is $s^{m_j}$ and its fibre multiplicity is exactly $m_j$. The map
\[
 J'_j\longrightarrow N_j\times_{D_j}\Delta_j,\qquad y\longmapsto(q_j(y),s_j(y))
\]
is finite, is an isomorphism over $s\ne0$, and has normal source. Locally at the central fibre the base-change ring has equation $z^{m_j}-s^{m_j}=0$. Its normalization is the disjoint union of branches $z=\zeta_j^rs$. The displayed map gives those branches in the covering charts. It is thus the normalization everywhere, with its specified root coordinate.

\par\smallskip\noindent\textbf{Step 2.} Put $\rho=e^{\pi i/3}$, $\mu_1=(2-\rho)/3$, $\mu_2=(1-i)/2$, $\kappa_1=\beta_c(z_1)+2\mu_1$, and $\kappa_2=\beta_c(z_2)+3\mu_2$. The fixed-point transformation laws for $\mu$ give the two displayed values directly. We have $\operatorname{Im}\kappa_j=D_c(z_j)<0$. Define
\[
 E_j=\mathbb C/(\mathbb Z+\kappa_j\mathbb Z),\quad
 F_1=\mathbb C/(\mathbb Z(\rho-2)-3\mathbb Z),\quad
 F_2=\mathbb C/(\mathbb Z(i-1)-2\mathbb Z).
\]
Let $e_1=(0,1,-2,0)^t,f_1=(0,0,-3,1)^t,e_2=(0,1,-1,0)^t,f_2=(0,0,-2,1)^t$. Direct solution of the fixed equations and inspection of the columns give
\[
 \Lambda^{A_1}=\mathbb Z\varepsilon+\mathbb Zd,\quad
 \Lambda^{A_2}=\mathbb Z\varepsilon'+\mathbb Zd,\quad
 (A_j-I)\Lambda=\mathbb Ze_j+\mathbb Zf_j.
\]
The image lattices are saturated: they are the integer solutions of $\gamma=0,2u_{\Lambda}+w_{\Lambda}+3\delta=0$ and $\gamma=0,u_{\Lambda}+w_{\Lambda}+2\delta=0$, respectively. The determinants of the fixed-plus-moving basis matrices are $-3,-2$.

\par\smallskip\noindent\textbf{Step 3.} The complex-linear maps $L_1(\xi,\eta)=(\eta,\xi-\eta/3)$ and $L_2(\xi,\eta)=(\eta,\xi-\eta/2)$ carry the product periods to those fixed and moving sublattices. They induce isogenies to the original central tori, with kernel generated respectively by $(2/3,\rho)$ and $(1/2,i)$, of exact orders three and two. For example $\widehat u=2d/3+e_1-2f_1/3=d/2+e_2-f_2/2$ gives both kernel representatives.

\par\smallskip\noindent\textbf{Step 4.} The affine deck maps lift to
\[
 (\xi,\eta)\longmapsto(\xi+\kappa_1/3,(\rho-1)\eta),\qquad
 (\xi,\eta)\longmapsto(\xi-\kappa_2/4,-i\eta).
\]
Substitution in $L_j$ gives exactly $R_jL_j+\Pi_cv_j/m_j$. They commute with their kernel translations: the moving differences are periods, since $(\rho-2)\rho=-1-\rho=-(\rho-2)-3$ can equally be checked in the lattice basis, and $(-i-1)i=1-i=-(i-1)$. More directly these assertions follow by the two-dimensional period matrices of the indicated rotations. A nonidentity element of either resulting product group translates the $E_j$ coordinate by $2a/3+b\kappa_1/3$ or $a/2-b\kappa_2/4$, respectively. Real independence of $1,\kappa_j$ proves that this translation vanishes only for the identity element. Hence the groups $C_3\times C_3$ and $C_2\times C_4$ act freely and have quotient precisely $S_j$.
\end{proof}
\begin{cited}[FAM A.2.1; K]\label{fam:A.2.1}
The rational centralizer of $(A_1,A_2)$ is $\mathbb QI+\mathbb QE$. Every biholomorphism $X_c\to X_{c'}$ has $c'-c=n\in2\mathbb Z$, complex-linear fibre part $I_2$, and source-to-target lattice map $I-nE$. In the full fixed lattices,
\[
 \Lambda^{A_1}/N_1\Lambda\simeq\mathbb Z/3,\qquad
 \Lambda^{A_2}/N_2\Lambda\simeq\mathbb Z/4\oplus\mathbb Z/2,
\]
with $d$ generating the second factor of the latter group and vanishing in the former.
\end{cited}
\uses{Definition~\ref{def:conventions}; FAM A.1.1 and preceding appendix calculations.}
\begin{proof}
\par\smallskip\noindent\textbf{Step 1.} The common invariant vector space is $\mathbb Qd$, and the common invariant covector space is $\mathbb Q\gamma$; this follows by intersecting the fixed bases in A.1.2 and the rows $\langle\gamma,2u_{\Lambda}+w_{\Lambda}+3\delta\rangle$, $\langle\gamma,u_{\Lambda}+w_{\Lambda}+2\delta\rangle$. A commuting endomorphism preserves $\mathbb Qd\subset\ker\gamma$. Its induced matrix on $\ker\gamma/\mathbb Qd$ commutes with
$Q_1=\left(\begin{smallmatrix}0&1\\-1&-1\end{smallmatrix}\right)$ and $Q_2=\left(\begin{smallmatrix}0&-1\\1&0\end{smallmatrix}\right)$. Commutation with $Q_2$ makes it $\left(\begin{smallmatrix}b&r\\-r&b\end{smallmatrix}\right)$, and the $Q_1$ equation sets $r=0$. Write the remaining columns as
$H\widehat\gamma=(a,x,y,k)^t$, $H\widehat u=b\widehat u+u_0d$, $H\widehat w=b\widehat w+w_0d$, $Hd=d_0d$. The equations on $\widehat u,\widehat w$ give $d_0-w_0=b+u_0$, $u_0=2w_0$ and $u_0=w_0$, hence $u_0=w_0=0,d_0=b$. The first-column equations give $x=2(a-b),y=-4(a-b),x=-y$. Thus $a=b,x=y=0$, proving the assertion. Conversely $A_jE=EA_j=E$.

\par\smallskip\noindent\textbf{Step 2.} A fibre biholomorphism lifts to an affine map on $\mathbb C^2$: the derivative of a lifted map is lattice periodic, descends to a holomorphic function on a compact torus, and is constant. Its lattice map is locally constant and commutes with both monodromies. Write it as $H=aI+bE$; integrality and $\det H=a^4=1$ give $a=\pm1,b\in\mathbb Z$. The equation $L\Pi_c=\Pi_{c'}H$, in columns three and four, gives $L=aI_2$; column one gives $a(c'-c)+b=0$. Thus $n=c'-c\in\mathbb Z$.

\par\smallskip\noindent\textbf{Step 3.} The normalized base-change construction in A.1.2 is functorial for biholomorphisms over the base: the base-changed isomorphism transports integral elements of the total quotient ring to integral elements, and its inverse does the same. Its lift is over the identity of the root disc. It commutes with the prescribed deck generator because both compositions lift the same map $(x,s)\mapsto(F(x),\zeta_js)$ on the unnormalized base change. This proves exact commutation, rather than commutation after an unchosen power of the generator.

\par\smallskip\noindent\textbf{Step 4.} Lift the translation of the resulting affine torus map holomorphically to a vector $B_j(s)$ on the simply connected root disc. Its deck equation, expressed in the source periods, is
\[
 B_j(g_jz)-R_j(z)B_j(z)=\Pi_c(g_jz)
 \left(\frac{(1-a)v_j+n\gamma(v_j)d}{m_j}+\ell_j\right),\qquad\ell_j\in\Lambda.
\]
The integral vector is constant by continuity and real invertibility of the period matrix. Iterating the equation $m_j$ times gives
\[
 N_j\ell_j=-(1-a)v_j-n\gamma(v_j)d.
\]
Exact multiplication yields
\[
 N_1=\begin{pmatrix}3&0&0&0\\6&0&0&0\\-12&0&0&0\\0&2&1&3\end{pmatrix},\quad
 N_2=\begin{pmatrix}4&0&0&0\\12&0&0&0\\-12&0&0&0\\0&2&2&4\end{pmatrix}.
\]
In fixed bases their images are $3\mathbb Z\varepsilon+\mathbb Zd$ and $4\mathbb Z\varepsilon'+2\mathbb Zd$. These explicit generators prove the asserted quotient groups. Applying $\gamma$ to the norm equation excludes $a=-1$ already at $m_1=3$, since $3\gamma(\ell_1)=-2$. For $a=1$, the second norm equation says exactly $2\mid n$; the first imposes no condition. Program \texttt{code/ab\_exact.py} verifies the matrices, Smith forms, centralizer and membership independently by Hermite-form comparisons. All its computations are over $\mathbb Z$ or $\mathbb Q$.
\end{proof}
\begin{cited}[FAM A.3.1; K]\label{fam:A.3.1}
For each $n\in2\mathbb Z$ there is a biholomorphism $F_{n,c}:X_c\to X_{c+n}$, holomorphic jointly in $c$, with $F_{2h,c+2k}F_{2k,c}=F_{2(h+k),c}$. The quotient family over $0<|w|<e^{-\pi M}$, $w=e^{-\pi ic}$, is proper and holomorphic, with pairwise non-biholomorphic fibres.
\end{cited}
\uses{Definition~\ref{def:conventions}; FAM A.1.1 and preceding appendix calculations.}
\begin{proof}
\par\smallskip\noindent\textbf{Step 1.} Integer shifts satisfy $\Pi_{c+n}=\Pi_c(I+nE)$, so the middle lattices agree. Set
\[
 r_1=-\tfrac n3(\widehat u+\widehat w),\quad
 r_2=\tfrac n4(\widehat u+\widehat w),\quad
 \ell_1=-n\widehat w,\quad \ell_2=\tfrac n2\widehat u,\quad b_j=\Pi_cr_j.
\]
The four-entry equations are $(I-A_1)r_1=(0,0,-n,n/3)^t$ and $(I-A_2)r_2=(0,n/2,0,-n/4)^t$. Consequently $B_j=b_j$ solves the preceding deck equation for $a=1$. The map $(z,\zeta)\mapsto(z,\zeta+b_j(z))$ extends on the entire root family and descends, together with its inverse, to each finite filling.

\par\smallskip\noindent\textbf{Step 2.} Its middle translation on the punctured overlap is
\[
 a_j=b_j+n\gamma(v_j)\frac{\log s_j}{2\pi i}e.
\]
This follows from the actual overlap $h_{j,c}(\zeta)=\zeta+\sigma_{j,c}$, where $\sigma_{j,c}=(\log s_j)/(2\pi i)\,\Pi_cv_j$: precisely $\sigma_{j,c+n}-\sigma_{j,c}=n\gamma(v_j)(\log s_j)/(2\pi i)e$. Changes of logarithm add periods. Since the chosen generator changes $(\log s_j)/(2\pi i)$ by $-1/m_j$, we have $a_j(g_jz)=R_ja_j(z)+\Pi_c(g_jz)\ell_j$.

\par\smallskip\noindent\textbf{Step 3.} On the universal cover of $B^\circ$, extend $\ell_1,\ell_2$ by the integral cocycle law $\ell_{gh}=\ell_g+A_g\ell_h$. Freeness of $\pi_1(B^\circ)$ makes this well defined. At $g_0=(g_1g_2)^{-1}$ it gives
\[
 \ell_0=-A_0(\ell_1+A_1\ell_2)=\tfrac{3n}2\widehat w-\tfrac n2d.
\]
Put $q_n=(-3n/2,n/2)^t$ and $a_0=sq_n$. Then $a_0(g_0z)-a_0(z)=\Pi_c\ell_0$. In cusp torus coordinates its map is $(x,t_c)\mapsto(t_c^{q_n}x,t_c)$. The integral linear map $(y,h)\mapsto(y+hq_n,h)$ of the fan lattice translates each height-one triangle by an integer vector. It therefore permutes every cone of the prescribed periodic triangulation. The map and its inverse extend across all strata, including both triple points and all three double curves after quotient. Moreover the cusp deck multipliers are identical at $c$ and $c+n$, since $C_{c+n}\lambda-C_c\lambda=(0,n\lambda_1)^t$ has integral exponential. Thus the extended map descends with all its nonconstant units retained.

\par\smallskip\noindent\textbf{Step 4.} We supply the global analytic correction. Quotient the trivial affine vector bundle on the universal cover by
$(z,y)\mapsto(gz,R_g(z)y+\Pi_c(gz)\ell_g)$. Its action law is exactly the cocycle law by period covariance. Specify at the finite points sections of the form $a_j+\eta_j$, with $\eta_j$ holomorphic equivariant on the full root disc; at the cusp specify $a_0+\eta_0(t_c)$ with $\eta_0$ holomorphic on the full disc. These prescriptions define a locally nonempty torsor under a rank-two locally free sheaf $\mathcal V$ on $B$.

\par\smallskip\noindent\textbf{Step 5.} Projection to the first coordinate gives
\[
 0\longrightarrow\mathcal O_Be\longrightarrow\mathcal V
 \longrightarrow\mathcal O_B(-p_0)\longrightarrow0.
\]
Indeed its factors are $-1/\tau,1/\tau,1$. Multiplication by $\Theta$ identifies this first-coordinate sheaf with holomorphic functions at ordinary and finite points and functions vanishing at $p_0$: the finite vanishing orders are $2,1$, and the cusp pole is one. At a finite stalk, any equivariant first component has an equivariant vector lift by averaging a preliminary lift $(h_{\rm loc},0)$ over the finite group. This proves stalkwise surjectivity and local freeness, including the off-diagonal entries of $R_j$.

\par\smallskip\noindent\textbf{Step 6.} Both $H^1(\mathbb P^1,\mathcal O)$ and $H^1(\mathbb P^1,\mathcal O(-1))$ vanish. To make the analytic gluing explicit, split a smooth cocycle with a partition of unity, solve $\bar\partial$ locally by the Cauchy integral, and compare the two resulting solutions on an annulus. A Laurent series $\sum a_nt^n$ is the difference of its nonnegative and negative parts for $\mathcal O$; for $\mathcal O(-1)$ use the second frame $t^{-1}$, so the negative terms become $t^{-1}(-\sum_{n<0}a_nt^{n+1})$. Both series converge on their respective discs and leave no exponent unsolved. This kills the cocycle for any refinement. The exact sequence gives $H^1(B,\mathcal V)=0$. Project the torsor cocycle to $\mathcal O(-1)$, solve it, lift the solutions, and solve the remaining scalar cocycle. We obtain one global section $a$, with $a=a_j+\eta_j$ at finite points and $a=sq_n+\eta_0$ at the cusp.

\par\smallskip\noindent\textbf{Step 7.} The three actual maps are now
\[
 \zeta\mapsto\zeta+a,\qquad
 \zeta\mapsto\zeta+b_j+\eta_j,\qquad
 (x,t_c)\mapsto(e^{2\pi i\eta_0(t_c)}t_c^{q_n}x,t_c).
\]
The finite overlap identity is the equality of sums
$\sigma_{j,c}+a=\sigma_{j,c+n}+b_j+\eta_j$. The cusp overlap identity is $e^{2\pi ia}=t_c^{q_n}e^{2\pi i\eta_0}$. The holomorphic-unit factors act on every toric chart. The local inverses satisfy the same identities; hence the maps glue to a biholomorphism.

\par\smallskip\noindent\textbf{Step 8.} For coherence, perform the additive solution only for $n=2$ and multiply that solution by $k$ for $n=2k$. The cocycles have $\gamma(\ell_g)=0$, so their affine terms $\Pi_c\ell_g$ are independent of $c$. The reference translations, $\mathcal V$, and the correction can all be chosen independent of $c$. Adding the translations and multiplying the toric factors proves the stated cocycle identity exactly.

\par\smallskip\noindent\textbf{Step 9.} Finally the original gluing constructs a holomorphic family over any parameter disc compactly contained in $\mathcal P$. In ordinary and finite charts, real invertibility of $\Pi_c$ is uniform on compact parameter sets, so a fixed compact real parallelepiped proves properness and the lattice quotient is holomorphic. At the cusp write $Z=sB_0+C(t_c,c)$ with $B_0=\left(\begin{smallmatrix}0&1\\-1&0\end{smallmatrix}\right)$. On a compact parameter disc $C$ is uniformly bounded; shrink the cusp radius until
$\|(-2\pi\operatorname{Im}C)/\log|t_c|\|<1/2$. The effective logarithmic lattice is then within $1/2$ of the unimodular $B_0$, uniformly. Reduce its logarithmic valuation modulo the effective logarithmic lattice to a bounded parallelogram. Choose a height-one triangle containing that reduced valuation. Its barycentric coordinates are nonnegative and sum to one; the corresponding three unimodular cone characters therefore have modulus at most one. This puts the point in the closed unit polydisc of that cone. Only finitely many triangles meet the fixed bounded parallelogram, uniformly on the parameter compactum. Boundary orbits have the same representatives by translating their height-one components to these triangles. Thus every orbit has a representative in this finite union of cone polydiscs. Those closed polydiscs times a compact parameter set form a compact surjective representative set. This proves properness at the cusp and proper discontinuity by the same uniform lower bound on nonzero lattice translations. The quotient charts are open subsets of a fixed smooth toric variety times the parameter disc, so projection to the parameter is a submersion. The explicit overlap maps are jointly holomorphic. These local families agree after shrinking cusp radii, and glue over $\mathcal P$. The coherent $2\mathbb Z$ action is properly discontinuous because its action on $\mathcal P$ is. Local inverse logarithms for $w=e^{-\pi ic}$ identify the quotient family with these proper local families. A.2.1 gives pairwise non-biholomorphic fibres.
\end{proof}
\begin{proposition}[FAM A.4.1; T]\label{fam:A.4.1}
With $\kappa_j(c)=\kappa_j^{\rm part}+c$, the Albanese elliptic curves of the reduced finite fibres are
\[
 \operatorname{Alb}(S_1)=\mathbb C/(\mathbb Z+\kappa_1(c)\mathbb Z),\qquad
 \operatorname{Alb}(S_2)=\mathbb C/(\mathbb Z+\tfrac12\kappa_2(c)\mathbb Z).
\]
Writing $J(q_\tau)=q_\tau^{-1}+744+196884q_\tau+\cdots$, their invariants are
\[
 j_1(w)=J(e^{-2\pi i\kappa_1^{\rm part}}w^2),\qquad
 j_2(w)=J(e^{-\pi i\kappa_2^{\rm part}}w).
\]
These are convergent formulas because the displayed nome moduli are below one. Outside the explicitly defined countable exceptional set in the proof, equality of both Albanese invariants forces $c'-c\in2\mathbb Z$. Thus the full generic Albanese cross-check of Lead 3 holds independently of the centralizer calculation.
\end{proposition}
\uses{Definition~\ref{def:conventions}; FAM A.1.1 and preceding appendix calculations.}
\begin{proof}
\par\smallskip\noindent\textbf{Step 1.} In the product covers of A.1.2 the invariant holomorphic one-forms are exactly the multiples of $d\xi$: the moving factor is multiplied by a primitive third or fourth root. Integration thus identifies the Albanese with the quotient of the fixed line by the fixed-coordinate projections of every deck translation. These are, respectively,
\[
 \mathbb Z+\kappa_1\mathbb Z+\tfrac23\mathbb Z+\tfrac{\kappa_1}3\mathbb Z
   =\tfrac13\mathbb Z+\tfrac{\kappa_1}3\mathbb Z,
\quad
 \mathbb Z+\kappa_2\mathbb Z+\tfrac12\mathbb Z+\tfrac{\kappa_2}4\mathbb Z
   =\tfrac12\mathbb Z+\tfrac{\kappa_2}4\mathbb Z.
\]
Multiplication by three and two gives the stated maps of complex tori. In particular the Albanese computation includes the kernel of the product isogeny and the affine quotient, rather than just its fixed eigenspace. Using the positive-imaginary representatives $-\kappa_1$ and $-\kappa_2/2$ in the standard modular $j$-series gives the formulas. At sufficiently small nome $J$ is injective: $1/J(q_\tau)=q_\tau+O(q_\tau^2)$ has nonzero derivative at zero. Thus $j_2$ alone distinguishes parameters within a sufficiently small punctured subdisc when both parameters lie there. For the full generic assertion put $x=-\kappa_2(c)/2$ and $d_*=\kappa_1^{\rm part}-\kappa_2^{\rm part}$, so $-\kappa_1(c)=2x-d_*$. The independently certified interval in \texttt{cd\_special\_intervals.py} gives $-0.2<\operatorname{Im}d_*<-0.1$; in particular it is nonzero. Equality of the two elliptic invariants at $c,c'$ gives matrices $M,N\in\operatorname{PSL}_2(\mathbb Z)$ with $x'=M(x)$ and
\[
 N(2x-d_*)=2M(x)-d_*.
\]
Define the exceptional set to be the $c\in\mathcal P$ obtained from roots of these equations for all pairs for which the two rational functions are not identical. Each equation has at most two roots after clearing denominators, so this is a countable, explicitly specified set. For a pair giving an identity, let $C(x)=2x-d_*$. Then $N=CMC^{-1}$. The real matrix $M$ preserves both the real circle $\mathbb R\cup\{\infty\}$ and $C^{-1}(\mathbb R\cup\{\infty\})$, the parallel line with imaginary part $\operatorname{Im}d_*/2$ together with infinity. Their unique intersection is infinity, so $M$ fixes infinity. An element of $\operatorname{PSL}_2(\mathbb Z)$ fixing infinity is $x\mapsto x+b$, $b\in\mathbb Z$, and then $N(y)=y+2b$. It follows that $c'-c=-2b$. This proves the generic necessity with its actual exceptional set; A.2.1 proves the necessity also on that set.
\end{proof}
\begin{cited}[FAM A.4.2; K]\label{fam:A.4.2}
For every $c\in\mathcal P_{\rm KS}$, the Kodaira--Spencer class of $\partial/\partial c$ in $H^1(X_c,T_{X_c})$ is nonzero. On any fixed ordinary fibre, put $T=\operatorname{Im}\tau$, $m=\operatorname{Im}\mu$, $a=m/T$, $D=\operatorname{Im}\beta_c-6m^2/T$. In the convention $\mathrm{KS}(\partial_c)=[\bar\partial h]$ for a smooth $(1,0)$ lift $h$, its restricted torus-deformation representative is
\[
 -\frac{1}{2iD}\,\partial_{\zeta_2}\otimes
       (d\bar\zeta_2-a\,d\bar\zeta_1).
\]
\end{cited}
\uses{Definition~\ref{def:conventions}; FAM A.1.1 and preceding appendix calculations.}
\begin{proof}
\par\smallskip\noindent\textbf{Step 1.} Retain the same real marking $r\in\mathbb R^4$ and set $\zeta=\Pi_cr$. The two imaginary-part equations give
$r_1=(\operatorname{Im}\zeta_2-a\operatorname{Im}\zeta_1)/D$. Thus at $c+\delta$ the marked real-linear identification is
\[
 \zeta^{\rm new}=\zeta+\frac{\delta}{2iD}P_{\rm KS}(\zeta-\bar\zeta),\qquad
 P_{\rm KS}=\begin{pmatrix}0&0\\-a&1\end{pmatrix},\qquad P_{\rm KS}^2=P_{\rm KS}.
\]
It is exactly $\Pi_{c+\delta}\Pi_c^{-1}$ as a real map, so it sends every marked lattice vector to the corresponding vector. Its determinant is $(D+\operatorname{Im}\delta)/D>0$. The transported $(0,1)$ graph coefficient is $\nu_\delta=\delta P_{\rm KS}/(2iD+\delta)$. The derivative at zero is $P_{\rm KS}/(2iD)$. Equivalently the smooth lift of $\partial_c$ in moving coordinates has vertical part $r_1\partial_{\zeta^{\rm new}_2}$ and $\bar\partial r_1=-(d\bar\zeta_2-a d\bar\zeta_1)/(2iD)$, giving the stated sign.

\par\smallskip\noindent\textbf{Step 2.} A nonzero constant tangent-valued $(0,1)$ form on a compact torus cannot be $\bar\partial$ of a global vector field. Write that field in the invariant holomorphic frame and integrate each derivative of its periodic coefficients against invariant probability volume. Every such derivative has integral zero, whereas the coefficient of $\partial_{\zeta_2}\otimes d\bar\zeta_2$ above has nonzero integral. This proves nonvanishing for the fibre deformation.

\par\smallskip\noindent\textbf{Step 3.} We explain its transfer to the threefold. Over the dual numbers $A=\mathbb C[\epsilon]/(\epsilon^2)$ at $c$, the chart proof of A.1.1 gives the relative line isomorphism $\omega^{-2}_{X_A/A}=f_A^*\mathcal O_B(p_2)$. Extra parameter differentials disappear in relative forms and all determinant characters are independent of $c$. The two pulled-back sections form an $A$-basis of $H^0(X_A,\omega^{-2})$: reduce a section modulo $\epsilon$, subtract its linear combination of the two lifts, and use the exact sequence $0\to L\xrightarrow{\epsilon}L_A\to L\to0$ to write the remainder uniquely as $\epsilon$ times another such combination. The same argument proves independence.

\par\smallskip\noindent\textbf{Step 4.} A first-order trivialization would preserve this complete system and give a base automorphism reducing to the identity. It fixes all three special values even over $A$. On the finite etale charts the relative critical ideal is $(s_j^{m_j-1})$, and $t_j=s_j^{m_j}$ is zero on that scheme; its constant parameter ring $A$ injects. At a cusp triple chart, the critical ideal of $t_c=z_0z_1z_2$ is $(z_0z_1,z_0z_2,z_1z_2)$ and has the same property. There are no other critical values. Therefore an infinitesimal displacement of a special value would remain nonzero in its critical chart, a contradiction. Writing a base infinitesimal automorphism as $t+\epsilon(b_0+b_1t+b_2t^2)$, fixing $0,1,\infty$ gives $b_0=b_1=b_2=0$. The trivialization would consequently trivialize the displayed nonzero torus deformation. Finally a zero Kodaira--Spencer class produces such a first-order trivialization by subtracting from local product charts the vector fields whose coboundary is the differentiated transition cocycle. This proves the claim.
\end{proof}
\subsubsection{Re-derivation of FAM A.5.1}
\begin{proof}
\par\smallskip\noindent\textbf{Step 1.} A.1.1 makes the base intrinsic. A.2.1 gives the necessary parity and excludes the negative linear part. A.3.1 supplies the inverse maps, all collar identities, and descent of the jointly holomorphic family. A.4.2 proves its infinitesimal effectiveness. Finally $|e^{-\pi ic}|=e^{\pi\operatorname{Im}c}$ and equality of exponentials is exactly difference in $2\mathbb Z$, so the coordinate has precisely the stated image and fibres.
\end{proof}
\subsubsection{The universal map statement}
The following source proof is that of \cite[Theorem 7.2]{DGKXv1}, with its parameter renamed; it is reproduced with that credit rather than attributed to the workbench.
\begin{proof}
\par\smallskip\noindent\textbf{Step 1.} Fix $c$ and let $F$ be a smooth torus fibre of $f_c$. Since $X_c$ is
diffeomorphic to $S^6$, $H^2(X_c,\mathbb Z)=0$. Thus
$c_1((h|_F)^*\mathcal O_{\mathbb P^1}(1))=0$. If $h|_F$ were nonconstant, the
pullback of the Fubini--Study form would be semipositive and nonzero
on an open subset of the K\"ahler surface $F$. Its wedge with a
K\"ahler form on $F$ would have strictly positive integral,
contradicting the vanishing of that Chern class. Hence $h$ is
constant on every smooth fibre.

\par\smallskip\noindent\textbf{Step 2.} The proper holomorphic map $(f_c,h):X_c\to\mathbb P^1\times\mathbb P^1$ has,
by Remmert's theorem, an irreducible compact analytic image $\Gamma$
of dimension one. Over the complement of the three special values,
$\Gamma$ is the graph of a holomorphic map. The first projection
$p:\Gamma\to\mathbb P^1$ is therefore generically one-to-one. It is also
finite: a positive-dimensional fibre of a map from an irreducible
compact analytic curve would be the whole curve, whereas $p$ is
surjective. Thus $p$ is finite and bimeromorphic. On a coordinate
disc of $\mathbb P^1$, the finite algebra $p_*\mathcal O_\Gamma$ sits in the
meromorphic function field of the disc and is integral over its
holomorphic functions. Normality of the disc forces equality, so
$p$ is an isomorphism. The second projection gives $\varphi$ and
$h=\varphi\circ f_c$. Surjectivity of $f_c$ gives uniqueness.
\end{proof}
